% =============================================================================
%  capitulos/cap6_testes.tex
%  CAPÍTULO 6 — TESTES E RESULTADOS
%
%  Objetivo: provar que o software funciona — ou documentar honestamente
%  o que não funcionou e por quê.
%
%  REGRA DE OURO: os resultados apresentados aqui devem se relacionar
%  diretamente com os objetivos definidos no Capítulo 1.
%  Ou seja: para cada objetivo específico, deve haver uma evidência aqui.
% =============================================================================

\chapter{Testes e Resultados}
\label{cap:testes}

TODO: Parágrafo introdutório descrevendo o contexto geral dos testes:
ambiente utilizado (local, servidor, nuvem), perfil dos usuários envolvidos
(se houver), e o que será apresentado nas próximas seções.


% -----------------------------------------------------------------------------
%  6.1 PLANO DE TESTES
%  Descreva o que foi testado, como e com qual objetivo.
%  Tipos de testes comuns em TCCs de software:
%    - Unitários: funções e métodos isolados
%    - Integração: comunicação entre módulos/camadas
%    - Usabilidade: experiência real com usuários
%    - Performance: tempo de resposta, carga
%    - Aceitação: validação com o cliente/orientador
% -----------------------------------------------------------------------------
\section{Plano de Testes}
\label{sec:plano-testes}

TODO: Descreva os tipos de testes realizados. Para cada tipo, mencione:
o que foi testado, qual framework ou abordagem foi usada, e qual era o
critério de sucesso.

\begin{itemize}
  \item \textbf{Testes Unitários}:
        TODO: Descreva quais módulos ou funções foram cobertos, qual
        framework foi utilizado (Jest, Pytest, JUnit, Vitest...) e qual
        foi a porcentagem de cobertura de código atingida.

  \item \textbf{Testes de Integração}:
        TODO: Descreva quais integrações foram validadas — por exemplo:
        API - banco de dados, frontend - backend. Mencione se foram
        usadas ferramentas como Postman, Insomnia, Supertest, etc.

  \item \textbf{Testes de Usabilidade}:
        TODO: Descreva como foram conduzidos (sessões com usuários,
        questionários SUS, heurísticas de Nielsen...), quem participou
        (perfil, quantidade), e quais métricas foram coletadas.
        O instrumento utilizado (questionário) pode ficar no Apêndice.

  \item \textbf{Testes de Performance} (se aplicável):
        TODO: Ferramentas utilizadas (Lighthouse, k6, JMeter...) e
        métricas avaliadas (tempo de resposta, número de requisições
        simultâneas, pontuação Lighthouse...).
\end{itemize}


% -----------------------------------------------------------------------------
%  6.2 RESULTADOS
%  Apresente os resultados de forma visual e objetiva.
%  Use prints das telas principais, gráficos de métricas e tabelas.
%  Organize por funcionalidade ou por objetivo específico.
% -----------------------------------------------------------------------------
\section{Resultados}
\label{sec:resultados}


% --- Telas principais do sistema ---
\subsection{Interface e Funcionalidades Principais}
\label{subsec:interface}

TODO: Apresente as capturas de tela mais representativas do sistema,
explicando o que cada tela mostra e como atende aos requisitos levantados.

\begin{figure}[H]
  \caption{TODO: Nome da tela — ex.: Tela Principal do Sistema}
  \label{fig:tela-principal}
  \centering
  % \includegraphics[width=0.9\textwidth]{imagens/tela_principal.png}
  \fbox{\rule{0pt}{5cm}\rule{0.85\textwidth}{0pt}} % ← Placeholder, remova depois
  \fonte{Elaborado pelo autor.}
\end{figure}

TODO: Explique o que a tela mostra e como atende ao(s) RF(s) correspondente(s).

% Repita o bloco acima para cada tela relevante:
\begin{figure}[H]
  \caption{TODO: Nome da tela 2}
  \label{fig:tela-2}
  \centering
  % \includegraphics[width=0.9\textwidth]{imagens/tela_2.png}
  \fbox{\rule{0pt}{5cm}\rule{0.85\textwidth}{0pt}} % ← Placeholder, remova depois
  \fonte{Elaborado pelo autor.}
\end{figure}


% --- Análise frente aos requisitos ---
\subsection{Análise Frente aos Requisitos}
\label{subsec:analise-requisitos}

TODO: Para cada objetivo específico definido no \cref{cap:introducao},
relate se foi atingido e apresente a evidência correspondente (tela,
teste, métrica).

A \cref{tab:resultados-rf} apresenta o status de implementação de cada
requisito funcional ao final do desenvolvimento.

\begin{table}[H]
  \caption{Status dos Requisitos Funcionais ao Final do Desenvolvimento}
  \label{tab:resultados-rf}
  \centering
  \begin{tabularx}{\textwidth}{llXl}
    \toprule
    \textbf{ID} & \textbf{Requisito} & \textbf{Evidência / Observação}
                & \textbf{Status} \\
    \midrule
    RF01 & TODO: Nome & TODO: Seção X / Figura Y demonstra o
           funcionamento...
         & \textcolor{green!50!black}{Implementado} \\
    RF02 & TODO: Nome & TODO: ...
         & \textcolor{green!50!black}{Implementado} \\
    RF03 & TODO: Nome & TODO: Implementado parcialmente devido a
           [motivo]...
         & \textcolor{orange}{Parcial} \\
    RF04 & TODO: Nome & TODO: Não implementado nesta versão devido
           a [motivo]...
         & \textcolor{red}{Não implementado} \\
    \bottomrule
  \end{tabularx}
  \fonte{Elaborado pelo autor.}
\end{table}


% --- Resultados dos Testes de Usabilidade (se aplicável) ---
% \subsection{Resultados dos Testes de Usabilidade}
% \label{subsec:usabilidade}
%
% TODO: Apresente os resultados do questionário ou sessão de usabilidade.
% Se usou o SUS (System Usability Scale), calcule a pontuação média e
% classifique conforme a escala (Excelente > 85, Bom > 72, Ok > 52...).
%
% \begin{table}[H]
%   \caption{Resultados do Questionário de Usabilidade (SUS)}
%   \label{tab:sus}
%   \centering
%   \begin{tabularx}{\textwidth}{lXc}
%     \toprule
%     \textbf{Participante} & \textbf{Perfil} & \textbf{Pontuação SUS} \\
%     \midrule
%     P1 & TODO: Perfil do participante & TODO: XX \\
%     P2 & TODO: Perfil do participante & TODO: XX \\
%     \midrule
%     \multicolumn{2}{l}{\textbf{Média}}    & TODO: XX \\
%     \bottomrule
%   \end{tabularx}
%   \fonte{Elaborado pelo autor.}
% \end{table}
